%%%%%%%%%%%%%%%%%%%%%%%%%%%%%%%%%%%%%%%%%%%%%%%%%%%%%
%			HLAVIČKA								%
%%%%%%%%%%%%%%%%%%%%%%%%%%%%%%%%%%%%%%%%%%%%%%%%%%%%%
\documentclass[openany,12pt]{memoir}

\usepackage[utf8]{inputenc}   
\usepackage[czech]{babel}
\usepackage[T1]{fontenc}
\usepackage[top=1.5cm, bottom=2cm, left=2cm, right=2cm]{geometry}  % --> NASTAVENÍ OKRAJŮ
\usepackage{fancyhdr}
\usepackage{graphicx}
\usepackage{lmodern}
\usepackage{xwatermark}
\usepackage{xcolor}



%%%%%% Package na zpěvník
\usepackage[full]{leadsheets}%http://mirrors.nic.cz/tex-archive/macros/latex/contrib/leadsheets/leadsheets_en.pdf   --> dokumentace	
\definesongtitletemplate{empty}{} 
\setchords{
format = \bfseries,   %tučné akordy
minor = {mi},% 
input-notation = {german},%
output-notation = {german}%
}
\definesongtitletemplate{empty}{} 

\newlength{\drop}
\newwatermark[allpages,color=red!50,angle=0,scale=2, xpos=0,ypos=0]{\includegraphics[width=5cm]{obr/pozadi.jpg}} %--> dvojka na pozadí


%%%% Vlastní příkazy
\newcounter{Slokočet}   %Automatické číslování slok
\newcommand{\mezera}{\vspace*{0.5cm}}   %Horizontální odsazení slok
\newcommand{\stred}{5.2cm}   %%% Na zarovnání slok doprostřed
\newcommand{\refren}{\mezera \noindent \textbf{R:} } %refrén
\newcommand{\sloka}{\mezera \noindent \addtocounter{Slokočet}{1} \arabic{Slokočet}. } 	%sloka, která se automaticky čísluje
\newcommand{\ssloka}{\mezera \noindent}  % vlastní číslo sloky


%%%%%%%%%%%%%%%%%%%%%%%%%%%%%%%%%%%%%%%%%%%%%%%%%%%%%%
%			NÁVOD									 %
%%%%%%%%%%%%%%%%%%%%%%%%%%%%%%%%%%%%%%%%%%%%%%%%%%%%%%
%1. Věci v hlavičce IGNOROVAT
%2. Píseň psát do prostoru mezi \begin{verse} a \end{verse}
%3. další řádek se značí dvěma odsazeníma (= dvakrát stisknout enter)
%4. \refren vždy na začátku refrenu a \sloka na začátek sloky (automaticky se čísluje)
%  \ = alt gr + q ; [/] = alt gr f/g ; {/} = alt gr + b/n; ^ = alt gr + 3 + mezera
%Cokoliv napíšete do ^{  } se bude brát jako akord
%když se toto bude dotýkat nějakého slova (nebude mezi tím a slovem mezera)
%tak se akordy zjeví nad slovem
%Když se to nedotýká slova, tak akord lítá ve vzduchu a vytiskne se větší mezera
%První možnost je asi preferovanější
%5. Akordy stačí psát jen do první sloky, když se nezmění -- kytaristi to zvládnou

%<++>
\usepackage{subfiles}
%</++>

\begin{document}
\pagestyle{plain}

%\documentclass[../main.tex]{subfiles}

\begin{song}{title=\centering Černej Pasažér \\\normalsize Traband \vspace*{-0.3cm}}  %% sem se napíše jméno songu a autor
\moveright 1cm \vbox{      %Varianta č. 1  ---> Jeden sloupec zarovnaný na střed
\begin{minipage}[t]{0.48\textwidth}\setlength{\parindent}{0.25cm}  %Varianta č. 2 --> Dva sloupce
\setcounter{Slokočet}{0}
\sloka
Mám ^{Dmi}kufr plnej přebytečnejch ^{A}krámů 
                        
a mapu zabalenou do ^{Dmi}plátna 
                              
Můj vlak však jede na opačnou ^{A}stranu 
                               
a moje jízdenka je dávno ^{Dmi}neplatná 

\mezera
\textbf{F  Dmi F  Dmi }

\sloka
Někde ve vzpomínkách stojí dům 

Ještě vidím, jak se kouří z komína 

V tom domě prostřený stůl 

Tam já a moje rodina 

\sloka
Moje minulost se na mě šklebí 

a srdce bolí, když si vzpomenu 

že stromy, který měly dorůst k nebi 

teď leží vyvrácený z kořenů 

\mezera
\textbf{F  Dmi F  Dmi }

\refren
Jsem černej^{B} pasažér 

^{C}nemám ^{F}cíl ani směr 

Vezu se ^{B}načerno ^{C}životem a ^{F}nevím 
             
Jsem černej^{B} pasažér 

^{C}Nemám ^{F}cíl ani směr 

Vezu se ^{B}odnikud ^{C}nikam a ^{A7}nevím, kde skončím

\end{minipage}\begin{minipage}[t]{0.48\textwidth}\setlength{\parindent}{0.45cm}  % V případě varianty č.2 jde odsud text do pravé části
\vspace*{0.6cm}

\sloka
Mám to všechno na barevný fotce  

někdy z minulýho století 

Tu jedinou a pocit bezdomovce 

si nesu s sebou jako prokletí

\mezera
\textbf{F  Dmi F  Dmi }

\refren

\sloka
Mám kufr plnej přebytečnejch krámů 

a mapu zabalenou do plátna 

Můj vlak však jede na opačnou stranu 

a moje jízdenka je dávno neplatná

\refren

\end{minipage}   %Součást druhé varianty
}
\end{song}
\setcounter{Slokočet}{0}
\newpage
%../main.tex]{subfiles}

\begin{song}{title=\centering Černobílý svět \\\normalsize Totální nasazení \vspace*{-0.3cm}}
%\vspace*{-1.2cm}
\moveright \stred \vbox{      %Varianta č. 1  ---> Jeden sloupec zarovnaný na střed
\setcounter{Slokočet}{0}
\sloka
^{Dmi}Jsme loutky ^{F}bez nití

^{B}já i všech ^{C{\color{white}\_\_\_}}patnáct druhů ^{Dmi}mých.

Bez smrti ^{F}bez žití

^{B{\color{white}\_\_\_\_}}bloumáme ^{C{\color{white}\_}}tupě po polích.

\refren
^{Dmi}Černá a ^{B{\color{white}\_}}bílá je,

^{C{\color{white}\_}}osm kroků tam a osm^{Dmi}sem,

frontová ^{B{\color{white}\_}}linie.

^{C}Na věky proti sobě ^{Dmi}jdem.

\sloka
^{Dmi}Stojí tu ^{F{\color{white}\_\_}}opodál

^{B}náš ^{{\color{white}\_\_\_}C}domnělý to ^{{\color{white}\_\_}Dmi}nepřítel.

I když jsem ^{F{\color{white}\_}}černý král

^{B{\color{white}\_\_}}nejsem svých ^{C\,\,\,\,}figur ^{{\color{white}\_}Dmi}velitel.

\sloka
Vždy první na tahu

jsem ten kdo bitvu rozpoutá,

ač nemá odvahu.

Bílý a nešťastný jsem král

\refren

\sloka
Přestal jsem počítat,

kolikrát jsem padnul a zas vstal,

kolikrát dostal mat,

útočil nebo utíkal

\sloka
Jsme jenom figury

ať král či pěšák beze jména

a hraní na vojáky je vaše lidská doména.

\refren

}
\end{song}
\setcounter{Slokočet}{0}
\newpage
\documentclass[../main.tex]{subfiles}

\chapter*{Do Ameriky}

\noindent\hspace{0.15\linewidth}
\begin{minipage}{0.7\linewidth}
\begin{song}{title=Do Ameriky}

\begin{verse}[numbered]
	^{D}Do Ameriky jezděj Parníky,\\
	když tam přijdeš, zdá se ti to ^{A7}všecko veliký,\\
	^{\phantom{S}}je to fakticky tuze praktický  \\
	když tam přijdeš, umět anglic^{D}ky. \\
\end{verse}

\begin{verse*}
	Refrain:
	^{D}Dobrou noc - good night, ^{A7}výborně - all right,\\
	Konrád White, už je off^{D}side (housemister - Dally sister!) \\
	^{D}Dneska je today, ^{A7}dobrý den - good day,\\
	já jsem byl taky chu^{D}dej. \\ 
	^{G}His Master's Voice, Yankee Doo^{G}dle,\\ \\
	^{En7}máš apetit? (mám) ^{A7}vem si štrůdl! (do pusy)\\
	^{D}Sto kilo je cent, ^{A7}patent je patent,\\
	husa v troubě happy^{D}end.
\end{verse*}

\begin{verse}[numbered]
	Co Buffalo Bill vrazil do kobyl,\\
	za to by si bejval byl koupil automobil,\\
	cowboy na koni už se nehoní,\\
	přestože mu Fordka nevoní.\\
\end{verse}

\begin{verse*}
	Refrain
\end{verse*}

\begin{verse}[numbered]
	Mám na západě chajdu v Nevadě,\\
	zlatou žílu v Arizoně, doly v Kanadě,\\
	babička na mě šetří v Panamě\\
	a já si tu žiju náramě!\\
\end{verse}



\end{song}

\end{minipage}\newpage

\end{document}
